\section{Discussion}
\label{section:disc}

The Raspberry Pi could detect the sensor but not read due to an I2C address connection. This was due to the small physical size of the sensor presented practical challenges during integration. With the sensor dimension of 2.5 x 3 x 0.83 mm and with limited laboratory equipment the group found it challenging, to manually solder connections with the required precision of 0.5 mm. This was something the group did not account for, and is something that future work will have to solve, either with larger sensor or more precise lab equipment that can assist with the soldering process required precision.

Reflecting on the test data created, one can see from the graphs that the filtered plot follows the gyro scope but with a small correction. This gave the expected result as the weight constant was so large. The code worked for test data and could read and calculate Pitch and Roll, but because of the mentioned hardware struggles, it could not be tested in real time. 

%Here you can discuss your results, limitations and new questions that have arose while doing the work. Depending on the size of this report, you can present discussion and conclusion in one common section or in two separate ones. 